%!TEX root = ../../main.tex
\section{같은 사례에서의 다른 결과 정의와 후속 결과 모형}
\label{sec:value-samecase}

서로 다른 결과 정의가 어떤 사례를 다르게 평가하고 어떤 후속 결과와 연결되는지를 같은 사례에서 본다. 이 절의 세 분석은 모두 보완 실험(15.16 노출 이후의 진단)이다.

\paragraph{교환가치 라벨과의 대응.} 킬·오브젝트·현상금에 고정 계수를 곱한 운영적 교환가치 라벨(v3.3 \texttt{market\_event}; 계수의 출처는 Baek과 Kwon \cite{baek2026killconditioned})을 $\SVI$와 같은 창 $(s,\ \eh + 1\,\text{ms})$에서 다시 계산하면(T \meN{} 행) 두 라벨의 부호는 \meFlip{}의 행에서 다르고 \meAgree{}에서 같다. 동결 $q$와 $\PTflex$의 예측을 고정한 채 교환가치 라벨로 점수화하면 $\dBrier$는 \meDbMarket{}으로 부호가 바뀐다($\SVI$ 라벨에서는 \meDbSvi{}). $q$는 $\SVI$ 라벨로 학습된 모델이므로 이 값은 표적 라벨에 대한 예측 비교의 의존성을 보여 주며, 두 라벨 가운데 어느 쪽이 더 나은 가치 지표인지에 관한 정보는 담지 않는다. 이 창에서 다시 계산한 교환가치 값은 학회 논문의 코퍼스와 창에서 얻은 값과 별개의 양이다.

\paragraph{후속 결과 모형.} 지표의 사후 활용을 보기 위해 결과 시점에서 평가하는 별도의 다항 로지스틱 모형 $R$을 두었다. 목표 $U$는 결과 시점 뒤 180초 안의 첫 엘리트 오브젝트 결과를 Blue, Red, 없음, 창 안 경기 종료, 동률로 나눈 것이다. $R_0$은 초기 승률·경기 시각·결과 시점·구간 길이·같은 구간의 물질 순차를 입력으로 하고, $R_1$은 여기에 $\dV$를 더한다. 두 모형은 TRAIN \rTrainN{} 행(out-of-fold $\dV$)에서 같은 예산으로 적합하고 Q\_SELECT \rQselN{} 행에서 $C = \rC$를 골랐으며, TEST \rTestN{} 행의 다항 Brier(클래스 합, 경기 가중)는 $R_0$ \rZeroBrier{}, $R_1$ \rOneBrier{}로 차이는 \rDelta{}이다. 귀속 팀이 Blue 또는 Red인 행으로 분모를 제한하면 \rZeroBR{}과 \rOneBR{}이다. $p_{\mathrm{post}} = \ppre + \dV$이므로 $R_1$은 결과 시점의 승률 정보를 추가한 모형과 재매개화 관계에 있다. 이 비교는 사전 예측기 $q$의 정확도와 구분되는 사후 활용 검사이며 구간을 계산하지 않았다.

\paragraph{사례 표본.} 킬 축과 $\SVI$가 갈리는 사례는 사전에 정한 셀(kill$+$ \& SVI$-$, kill$-$ \& SVI$+$, kill $= 0$)에서 sha256 순서로 셀당 8건까지 뽑아 기록하였다(부록 \ref{app:repro}). 질의 시각에서 사건 접두, $\Vhat$ 점수, $\SVI$, $q$까지 이어지는 사례별 추적 덤프는 미완료이다(부록 \ref{app:deferred}).
